\documentclass[11pt]{article}
\usepackage{amsmath,amssymb,amsthm,mathtools}
\usepackage[margin=1in]{geometry}
\theoremstyle{plain}
\newtheorem{theorem}{Theorem}
\newtheorem{lemma}{Lemma}
\theoremstyle{remark}
\newtheorem{remark}{Remark}
\newcommand{\E}{\mathbb{E}}\newcommand{\R}{\mathbb{R}}\newcommand{\Z}{\mathbb{Z}}
\newcommand{\ee}{\mathcal E}\newcommand{\norm}[1]{\lVert #1\rVert}

\title{Formalization brief: \texttt{lem\_free}\\
(on-diagonal heat-kernel bound via Nash's method)}
\date{6 July 2026}

\begin{document}
\maketitle

\noindent\textbf{Goal.} Remove the \texttt{sorry} at \texttt{lem\_free} in
\texttt{TypeDDecouplingLCLT.lean} (attached): for
$W$ a \texttt{DriftlessReversibleWalk} (kernel $p$, rates, reversing measure
$c_1\le m\le c_2$, exit rate $\le\Lambda$, range $\le\varrho$, nearest-neighbour
conductance $m(x)\,\mathrm{rate}(x,x{+}1)\ge\delta>0$),
\[
  \exists\,C>0:\quad \forall t\ge0,\ \forall r,\qquad
  p_t(r)\ \le\ \frac C{\sqrt{1+t}} .
\]
This is the paper's own Nash/CKS argument (Lemma 4.1). You may modify
\texttt{TypeDDecouplingLCLT.lean} only by replacing this \texttt{sorry} with a
proof and adding supporting lemmas (a separate file
\texttt{TypeDDecouplingNash.lean} for the abstract analysis is encouraged ---
keep it library-clean). Do not alter the statement or the structure
\texttt{DriftlessReversibleWalk}; do not touch other leaves. Whole project must
build; standard axioms only.

\section*{Tier 1: the discrete Nash inequality on $\Z$ (elementary, no Fourier)}

For finitely supported $f:\Z\to\R$, with
$\norm{\nabla f}_2^2:=\sum_x(f(x{+}1)-f(x))^2$:
\begin{itemize}
\item[(a)] \underline{1D Agmon bound:}
$\norm f_\infty^2\le 2\,\norm f_2\,\norm{\nabla f}_2$.
\emph{Proof:} $f(x)^2=\sum_{y\le x}\big(f(y)^2-f(y{-}1)^2\big)
=\sum_{y\le x}(f(y)+f(y{-}1))(f(y)-f(y{-}1))
\le2\norm f_2\norm{\nabla f}_2$ by Cauchy--Schwarz.
\item[(b)] \underline{Nash:} $\norm f_2^2\le\norm f_\infty\norm f_1$, so by (a),
$\norm f_2^4\le\norm f_1^2\cdot2\norm f_2\norm{\nabla f}_2$, i.e.
\[
  \norm f_2^{6}\ \le\ 4\,\norm f_1^{4}\,\norm{\nabla f}_2^{2}.
\]
\end{itemize}
Extend from finite support to $\ell^1\cap\ell^2$ by truncation/monotone limits
as needed (the application is to kernels, which are summable).

\section*{Tier 2: energy identity and Nash iteration}

With the Dirichlet form
$\ee(f):=\tfrac12\sum_{x,y}m(x)\,\mathrm{rate}(x,y)\,(f(y)-f(x))^2$:
\begin{itemize}
\item[(c)] \underline{Conductance lower bound:}
$\ee(f)\ge\tfrac12\sum_x m(x)\mathrm{rate}(x,x{+}1)(f(x{+}1)-f(x))^2
\ge\tfrac\delta2\norm{\nabla f}_2^2$ (drop all non-nearest-neighbour terms;
they are nonnegative).
\item[(d)] \underline{Energy identity:} for $u(t):=\sum_x m(x)\,p_t(x)^2$
(the $\ell^2(m)$ norm of the kernel started at the fixed origin --- adapt to
however \texttt{IsTransitionKernel} presents the forward/backward equations),
$u'(t)=-2\,\ee(p_t)$, from the Kolmogorov equation and reversibility.
\underline{Sanctioned fallback:} if the abstract kernel structure genuinely does
not provide enough regularity to differentiate under the sum, this single
identity (or its integrated form
$u(t)+2\int_s^t\ee(p_\tau)\,d\tau=u(s)$) may enter as ONE named hypothesis on
the structure; everything else must be proved.
\item[(e)] \underline{Iteration:} $\norm{p_t}_{\ell^1(m)}$ is conserved/bounded
($\le c_2/c_1$ appropriately normalised); Tier 1 + (c) give
$u'\le-\kappa\,u^3$ with $\kappa=\kappa(c_1,c_2,\delta)>0$; the ODE comparison
(for a nonincreasing $u\ge0$ with $u'\le-\kappa u^3$, one has
$u(t)\le(2\kappa t)^{-1/2}$ --- integrate $-u'/u^3\ge\kappa$) yields
$u(t)\le C/\sqrt t$.
\end{itemize}

\section*{Tier 3: from on-diagonal to the uniform bound}

\begin{itemize}
\item[(f)] By Cauchy--Schwarz and the semigroup property (Chapman--Kolmogorov),
$p_{2t}(r)\ \le$ (a product of two time-$t$ $\ell^2(m)$ norms, weighted by the
$m$-bounds) $\le\ C'\,u(t)^{1/2}\cdot(\text{const})\ \le\ C''/\sqrt t$ ---
carry the $c_1,c_2$ factors explicitly.
\item[(g)] For $t\le1$: $p_t(r)\le1$ (probabilities), and
$1\le\sqrt2/\sqrt{1+t}$. Combine (f)--(g) into the stated
$\exists C,\ \forall t\ge0:\ p_t(r)\le C/\sqrt{1+t}$.
\end{itemize}

\begin{remark}[hygiene and fallbacks]
(1) Search Mathlib first for any usable pieces (discrete Sobolev/Nash
inequalities, ODE comparison lemmas, Gronwall variants); reuse what exists.
(2) The only sanctioned named hypothesis is the energy identity (d), and only
if the kernel structure forces it; the Nash inequality (Tier 1), the
iteration (e), and the off-diagonal step (f) must be proved outright.
(3) Numerical constants are free --- any explicit $C(c_1,c_2,\delta,\Lambda,\varrho)$
suffices; do not chase optimality.
(4) Mind $t=0$ and the degenerate cases ($u$ may be $+\infty$ if the kernel
were not square-summable --- it is, being bounded by $1$ and summable; record
this).
(5) Report in the summary whether (d) was proved or hypothesised.
\end{remark}

\end{document}
